\documentclass{book}

\title{Analisi Matematica T-A}
\author{Michele Mazza}
\date{November 2023}


% Pacchetti di Codifica e Lingua
\usepackage[utf8]{inputenc}
\usepackage[T1]{fontenc}
\usepackage[italian]{babel}

% Pacchetti Grafici e Diagrammi
\usepackage{graphicx}
\usepackage{tikz}
\usepackage{tikz-3dplot}
\usetikzlibrary{angles,quotes,arrows.meta, positioning}
\usepackage{pgfplots}
\pgfplotsset{compat=1.18}

% Pacchetti Matematici
\usepackage{amsmath, amssymb, amsthm}
\usepackage{polynom}

% Pacchetti per Tabelle e Unità di Misura
\usepackage{geometry}
\usepackage{siunitx}
\usepackage{booktabs}
\usepackage{tabularx}
\usepackage{colortbl}

% Pacchetti per Layout e Formattazione
\usepackage{fancyhdr}
\usepackage{float}
\usepackage{microtype}

% Altri Pacchetti Utili
\usepackage{mdframed}
\usepackage{animate}

% Definizioni dei Colori
\definecolor{lightblue}{HTML}{9cd8fc}
\definecolor{myblue}{HTML}{139FE9}

% Definizione del comando \norm
\newcommand{\norm}[1]{\left\lVert#1\right\rVert}

% Ambienti Personalizzati
\newmdtheoremenv[
        backgroundcolor=orange!10,
        linecolor=myblue,
        linewidth=0pt,
        leftmargin=0pt,
        rightmargin=0pt,
        innerleftmargin=10pt,
        innerrightmargin=10pt,
        innertopmargin=10pt,
        innerbottommargin=10pt
        nobreak
]{curiosity}{Curiosità}

\newmdtheoremenv[
        font=\normalfont\small,
        linewidth=0pt,
        innerleftmargin=10pt,
        innerrightmargin=10pt,
]{example}{Esempio}

\newmdtheoremenv[
    backgroundcolor=lightblue,
    linecolor=myblue,
    linewidth=1.5pt,
    linewidth=0pt,
    leftmargin=0pt,
    innertopmargin=10pt,
    innerbottommargin=10pt,
    innerrightmargin=10pt,
    innerleftmargin=10pt
    nobreak
]{theorem}{Teorema}

\newmdtheoremenv[
    backgroundcolor=gray!10,
    linewidth=0pt,
    linewidth=0pt,
    leftmargin=0pt,
    innertopmargin=10pt,
    innerbottommargin=10pt,
    innerrightmargin=10pt,
    innerleftmargin=10pt
    nobreak
]{definition}{Definizione}

\newmdtheoremenv[
    backgroundcolor=gray!10,
    linecolor=myblue,
    linewidth=1.5pt,
    linewidth=0pt,
    leftmargin=0pt,
    innertopmargin=10pt,
    innerbottommargin=10pt,
    innerrightmargin=10pt,
    innerleftmargin=10pt
    nobreak
]{proposition}{Proposizione}

% Configurazioni di fancyhdr
\pagestyle{fancy}
\fancyhf{} % Pulisce tutti i campi
\fancyhead[LE,RO]{\thepage} % Numero di pagina
\fancyhead[RE]{\nouppercase{\leftmark}} % Titolo del capitolo
\fancyhead[LO]{\nouppercase{\rightmark}} % Titolo della sezione
\renewcommand{\headrulewidth}{0.5pt}
\renewcommand{\footrulewidth}{0pt}

\geometry{a4paper, left=20mm, right=20mm, top=20mm, bottom=20mm}

% Pacchetto Hyperref (inserito per ultimo)
\usepackage{hyperref}

% Additional Packages
\usepackage{xparse,array}

\ExplSyntaxOn

\NewDocumentCommand{\ruffini}{mmmm}
 {% #1 = polynomial, #2 = divisor, #3 = middle row, #4 = result
    \franklin_ruffini:nnnn { #1 } { #2 } { #3 } { #4 }
 }

\seq_new:N \l_franklin_temp_seq
\tl_new:N \l_franklin_scheme_tl
\int_new:N \l_franklin_degree_int

\cs_new_protected:Npn \franklin_ruffini:nnnn #1 #2 #3 #4
 {
    % Start the first row
    \tl_set:Nn \l_franklin_scheme_tl { & }
    % Split the list of coefficients
    \seq_set_split:Nnn \l_franklin_temp_seq { , } { #1 }
    % Remember the number of columns
    \int_set:Nn \l_franklin_degree_int { \seq_count:N \l_franklin_temp_seq }
    % Fill the first row
    \tl_put_right:Nx \l_franklin_scheme_tl
     { \seq_use:Nn \l_franklin_temp_seq { & } }
    % End the first row and leave two empty places in the next
    \tl_put_right:Nn \l_franklin_scheme_tl { \\ #2 & & }
    % Split the list of coefficients and fill the second row
    \seq_set_split:Nnn \l_franklin_temp_seq { , } { #3 }
    \tl_put_right:Nx \l_franklin_scheme_tl
     { \seq_use:Nn \l_franklin_temp_seq { & } }
    % End the second row
    \tl_put_right:Nn \l_franklin_scheme_tl { \\ \hline }
    % Split and fill the third row
    \seq_set_split:Nnn \l_franklin_temp_seq { , } { #4 }
    \tl_put_right:Nx \l_franklin_scheme_tl
     { & \seq_use:Nn \l_franklin_temp_seq { & } }
    % Start the array (with \use:x because the array package
    % doesn't expand the argument)
    \use:x
     {
        \exp_not:n { \begin{array} } { r | *{\int_eval:n { \l_franklin_degree_int - 1 }} { r } | r }
     }
    % Body of the array and finish
    \tl_use:N \l_franklin_scheme_tl
    \end{array}
 }
\ExplSyntaxOff

\begin{document}

\frontmatter
\maketitle
\tableofcontents

\chapter*{Prefazione}
Il presente libro rappresenta un riassunto del corso di Analisi Matematica T-A tenuto dal Prof.re Cataldo Grammatico presso l'Università di Bologna. Rivolto agli studenti del corso di laurea in Ingegneria Gestionale (codice laurea 0925, Tipo di laurea L-9, Codice corso 28622).

\mainmatter

\chapter{Elementi di base}

\section{Insiemi}
\begin{center}
    \begin{tabular}{|c|c|c|}
        \hline
        \rowcolor{lightblue}
        Nome                & Simbolo            & Descrizione                          \\
        \hline
        Insieme vuoto       & \(\emptyset\)      & Insieme senza elementi               \\
        Unione              & \(A \cup B\)       & Elementi in A o in B o in entrambi   \\
        Intersezione        & \(A \cap B\)       & Elementi comuni a A e B              \\
        Differenza          & \(A \setminus B\)  & Elementi in A ma non in B            \\
        Complementare       & \(A^c\)            & Tutti gli elementi non in A          \\
        Prodotto Cartesiano & \(A \times B\)     & Coppie ordinate di elementi da A e B \\
        Potenza             & \(\mathcal{P}(A)\) & Insieme di tutti i sottoinsiemi di A \\
        \hline
    \end{tabular}
\end{center}

\section{Logica Matematica}
\begin{center}
    \begin{tabular}{|c|c|c|}
        \hline
        \rowcolor{lightblue}
        Simbolo             & Nome         & Significato                 \\
        \hline
        \(\land\)           & AND          & E logico                    \\
        \(\lor\)            & OR           & O logico                    \\
        \(\neg\)            & NOT          & Negazione                   \\
        \(\Rightarrow\)     & Implica      & Implicazione                \\
        \(\Leftrightarrow\) & Se e solo se & Equivalenza                 \\
        $\forall$           & Per ogni     & Quantificatore universale   \\
        $\exists$           & Esiste       & Quantificatore esistenziale \\
        \hline
    \end{tabular}
\end{center}

\section{Insiemi Numerici}
\begin{center}
    \begin{tabular}{|c|c|c|c|}
        \hline
        \rowcolor{lightblue}
        Insieme        & Nome      & Proprietà           & Esempi                        \\
        \hline
        \(\mathbb{N}\) & Naturali  & Solo positivi       & 1, 2, 3                       \\
        \(\mathbb{Z}\) & Interi    & Positivi e negativi & -2, -1, 0, 1                  \\
        \(\mathbb{Q}\) & Razionali & Rapporto di interi  & \(\frac{1}{2}, \frac{-3}{4}\) \\
        \(\mathbb{R}\) & Reali     & Include irrazionali & \(\sqrt{2}, \pi\)             \\
        \(\mathbb{C}\) & Complessi & Include immaginari  & \(1 + 2i, -i\)                \\
        \hline
    \end{tabular}
\end{center}

\section{Metodi di Dimostrazione in Matematica}
\noindent\begin{minipage}{\textwidth}
    \begin{theorem}
        \subsection{Principio di Induzione}
        Sia \( P(n) \) una proprietà definita per ogni intero \( n \geq n_0 \), dove \( n_0  \geq 0\) è un intero fissato. Se
        \begin{enumerate}
            \item \( P(n_0) \) è vera (caso base),
            \item per ogni \( n \geq n_0 \), l'implicazione \( P(n) \Rightarrow P(n+1) \) è vera (passo induttivo),
        \end{enumerate}
        allora la proprietà \( P(n) \) è vera per ogni intero \( n \geq n_0 \).

        \textbf{Esempio:}
        \(P(n)\): La somma degli angoli interni di un poligono di \( n \) lati è \( (n-2)\pi \) radianti.
        \begin{enumerate}
            \item \textbf{Passo Base:} Consideriamo un triangolo (un poligono con \( n = 3 \) lati). La somma dei suoi angoli interni è \( \pi \) radianti. Poiché \( (3-2)\pi = \pi \), il teorema è valido per \( n = 3 \).
            \item \textbf{Passo Induttivo:} Supponiamo che \( P(n) \) sia vera per un poligono di \( n \) lati. Consideriamo un poligono di \( n+1 \) lati. Tracciamo una diagonale che divide il poligono in un triangolo e un poligono di \( n \) lati. La somma degli angoli interni del triangolo è \( \pi \) radianti, e per l'ipotesi induttiva, la somma degli angoli interni del poligono di \( n \) lati è \( (n-2)\pi \) radianti. Quindi, la somma degli angoli interni del poligono di \( n+1 \) lati è \( \pi + (n-2)\pi = (n-1)\pi \), che corrisponde esattamente a \( (n+1-2)\pi \). Pertanto, la proprietà è vera per un poligono di \( n+1 \) lati.
        \end{enumerate}
    \end{theorem}
\end{minipage}
\begin{figure}[H]
    \centering
    \begin{tikzpicture}
        \filldraw[fill=lightblue, draw=myblue, thick] (0,0) -- (2,0) -- (2.5,1.5) -- (1.5,3) -- (0,2.5) -- cycle;
        \draw[dashed] (0,0) -- (1.5,3);
        \node at (3,1) {\( n \)-gon};
        \node at (-1,2) {Triangle};
    \end{tikzpicture}
    \caption{Poligono di n-lati}
\end{figure}

\begin{theorem}
    \subsection{Disuguaglianza di Bernoulli}
    Per ogni numero reale \( r \geq -1 \) e per ogni numero intero \( n \geq 0 \), vale la seguente disuguaglianza:
    \[ (1 + r)^n \geq 1 + nr \]

    \textbf{Osservazioni:}
    \begin{enumerate}
        \item L'elemento \(1\) rappresenta il valore iniziale, o il capitale iniziale in contesti finanziari, prima di qualsiasi aumento.
        \item L'espressione \( (1 + r) \) rappresenta un aumento del valore iniziale di una frazione \( r \). Per esempio, \( r = 0.05 \) significa un incremento del 5% rispetto al valore iniziale.
        \item \( (1 + r)^n \) descrive la crescita composta per \( n \) periodi, dove non solo l'incremento \( r \), ma l'intero valore \( 1 + r \), viene moltiplicato ripetutamente.
        \item Al contrario, \( 1 + nr \) rappresenta una crescita lineare, dove solo l'incremento \( r \) viene sommato ripetutamente al valore iniziale. Qui, \( 1 \) rimane costante, mentre \( nr \) è la somma totale degli incrementi.
        \item La disuguaglianza illustra che la crescita composta (elevamento a potenza) è più efficace della semplice somma ripetuta di un incremento, specialmente per valori grandi di \( n \) e per \( r \) positivo.
    \end{enumerate}
\end{theorem}

\begin{figure}[H]
    \centering
    \begin{tikzpicture}
        \begin{axis}[
                title={Confronto tra Crescita Esponenziale e Lineare},
                xlabel={$n$},
                ylabel={Valore},
                xmin=0, xmax=10,
                ymin=0, ymax=20,
                grid=both,
                axis lines=middle,
                axis line style=->,
                grid style={line width=.1pt, draw=gray!10},
                major grid style={line width=.2pt,draw=gray!50},
                minor tick num=5,
                enlargelimits={abs=0.5},
                ticklabel style={font=\small},
                xlabel style={at={(ticklabel* cs:1)}, anchor=north west},
                ylabel style={at={(ticklabel* cs:1)}, anchor=south west},
                legend pos=south east
            ]
            \addplot+[myblue, thick, smooth] plot coordinates {
                    (0,1)(1,2)(2,4)(3,8)(4,16)(5,32)
                };
            \addlegendentry{$(1+r)^n$}
            \addplot+[red, thick, smooth] plot coordinates {
                    (0,1)(1,2)(2,3)(3,4)(4,5)(5,6)(6,7)(7,8)(8,9)(9,10)(10,11)
                };
            \addlegendentry{$1+nr$}
        \end{axis}
    \end{tikzpicture}
    \caption{Confronto tra Crescita Esponenziale e Lineare per \( r = 1 \)}
\end{figure}

\section{Valore Assoluto}
\begin{definition}
    \subsection{Valore Assoluto (o modulo)}
    Il valore assoluto di un numero reale \( x \) è definito come
    \[ |x| =
        \begin{cases}
            x  & \text{se } x \geq 0, \\
            -x & \text{se } x < 0.
        \end{cases} \]
    Questa funzione restituisce la distanza di \( x \) dall'origine sul numero reale.
\end{definition}

\begin{theorem}
    \subsection{Disuguaglianza Triangolare, (Cauchy, Shwarz)}
    Per tutti i numeri reali \( x \) e \( y \), vale la seguente disuguaglianza:
    \[ |x + y| \leq |x| + |y| \]
    Questa disuguaglianza afferma che la magnitudine della somma di due numeri reali non è mai maggiore della somma delle loro magnitudini.
\end{theorem}

\begin{figure}[H]
    \centering
    \begin{tikzpicture}
        % Triangolo
        \coordinate (A) at (0,0);
        \coordinate (B) at (1,2);
        \coordinate (C) at (3,0);

        % Riempi il triangolo di un colore
        \fill[lightblue] (A) -- (B) -- (C) -- cycle;

        % Lati
        \draw [->,thick](A) -- (B);
        \node[left] at (B) {$x$};
        \draw [->,thick](B) -- (C);
        \node[right] at (C) {$y$};
        \draw [<-,thick](C) -- (A);
        \node[below] at (A) {$z = |x+y|$};
    \end{tikzpicture}
    \caption{Magnitudine della somma di due numeri reali}
\end{figure}

\begin{figure}[H]
    \centering
    \begin{tikzpicture}
        \coordinate (A') at(0,-1);
        \coordinate (B') at(2.23,-1);
        \coordinate (C') at(4.472,-1);
        \draw [->,thick](A') -- (B');
        \draw [->,thick](B') -- (C');
        \node at (2.23, -1.5) {$|x|+|y|$};
    \end{tikzpicture}
    \caption{Somma delle magnitudini di due numeri reali}
\end{figure}

\begin{theorem}
    \subsection{Proprietà del Prodotto}
    Per tutti i numeri reali \( x \) e \( y \), vale la seguente proprietà:
    \[ |xy| = |x||y| \]
    Questa proprietà afferma che il valore assoluto del prodotto di due numeri reali è uguale al prodotto dei loro valori assoluti.
\end{theorem}

\section{Intervalli}
\begin{definition}
    \subsection{Intervallo}
    Un \textbf{intervallo} è insieme di numeri tra due estremi, \(a\) e \(b\).
    \begin{itemize}
        \item \text{Chiuso} \([a, b]\)
        \item \text{Aperto} \((a, b)\)
    \end{itemize}
\end{definition}
\begin{figure}[H]
    \centering
    \begin{tikzpicture}
        \draw[-latex] (-1,0) -- (5,0) ; % Linea numerica
        \foreach \x in  {1,3} % Punti
        \draw[shift={(\x,0)},color=black] (0pt,2pt) -- (0pt,-2pt);
        \filldraw [myblue] (1,0) circle (2pt); % Cerchio pieno per a
        \filldraw [myblue] (3,0) circle (2pt); % Cerchio pieno per b
        \draw[thick, myblue](1,0) -- (3,0);
        \draw (1,0);
        \node[below] at (1,0) {$a$};
        \draw (3,0);
        \node[below] at (3,0) {$b$};
        \draw (5,-0.5);
        \node[right] at (5,-0.5) {\([a, b]\)};
    \end{tikzpicture}
    \caption{Intervallo aperto}
\end{figure}

\begin{figure}[H]
    \centering
    \begin{tikzpicture}
        \draw[->] (-1,-1) -- (5,-1) ; % Linea numerica
        \draw [myblue](1,-1) circle (2pt); % Cerchio vuoto per a
        \draw [myblue](3,-1) circle (2pt); % Cerchio vuoto per b
        \draw[thick, myblue](1.05,-1) -- (2.95,-1);
        \draw (1,-1);
        \node[below] at (1,-1) {$a$};
        \draw (3,-1);
        \node[below] at (3,-1) {$b$};
        \draw (5,-1.5);
        \node[right] at (5,-1.5) {\((a, b)\)};
    \end{tikzpicture}
    \caption{Intervallo chiuso}
\end{figure}
\begin{figure}[H]
    \centering
    \begin{tikzpicture}
        \draw[->] (-1,-1) -- (5,-1) ; % Linea numerica
        \filldraw [myblue](1,-1) circle (2pt); % Cerchio vuoto per a
        \draw[thick, myblue](1.05,-1) -- (5,-1);
        \draw (1,-1);
        \node[below] at (1,-1) {$a$};
        \draw (5,-1);
        \node[below] at (5,-1) {$\infty$};
        \draw (5,-1.5);
        \node[right] at (5,-1.5) {\([a, +\infty)\)};
    \end{tikzpicture}
    \caption{Intervallo semiaperto a sinistra}
\end{figure}

\begin{definition}
    \subsection{Maggiorante e Minorante}
    Sia A un intervallo, (chiuso o aperto).
    \begin{itemize}
        \item M è un \textbf{maggiorante} di A se \[\forall x \in A, \, x \leq M.\]
        \item m è un \textbf{minorante} di A se \[\forall x \in A, \, x \geq m.\]
    \end{itemize}
\end{definition}

\begin{figure}[H]
    \centering
    \begin{tikzpicture}
        \draw[->] (-1,0) -- (5,0) ; % Linea numerica
        \filldraw [myblue] (1,0) circle (2pt); % Cerchio pieno per a
        \draw [myblue] (3,0) circle (2pt); % Cerchio vuoto per b
        \filldraw [black] (0.4,0) circle (1pt); % Cerchio pieno per a
        \filldraw [black] (3.6,0) circle (1pt); % Cerchio vuoto per b
        \draw[thick, myblue](1,0) -- (3,0);
        % Aggiungi le parentesi e le etichette
        \node at (0.4,0.5) [above] {$m$};
        \node at (3.6,0.5) [above] {$M$};
        \node at (1,-0.5) [below] {$a$};
        \node at (3,-0.5) [below] {$b$};
        \node at (5,-0.5) [below] {\([a, b)\)};
    \end{tikzpicture}
    \caption{minorante e maggiorante di A}
\end{figure}
\begin{figure}[H]
    \centering
    \begin{tikzpicture}
        \draw[->] (-1,0) -- (5,0) ; % Linea numerica
        \filldraw [myblue] (1,0) circle (2pt); % Cerchio pieno per a
        \draw [myblue] (3,0) circle (2pt); % Cerchio vuoto per b
        \filldraw [black] (1,0) circle (1pt); % Cerchio pieno per a
        \filldraw [black] (3,0) circle (1pt); % Cerchio vuoto per b
        \draw[thick, myblue](1,0) -- (3,0);
        % Aggiungi le parentesi e le etichette
        \node at (1,0.5) [above] {$m$};
        \node at (3,0.5) [above] {$M$};
        \node at (1,-0.5) [below] {$a$};
        \node at (3,-0.5) [below] {$b$};
        \node at (5,-0.5) [below] {\([a, b)\)};
    \end{tikzpicture}
    \caption{minorante e maggiorante di A che coincidono con gli estremi dell'insieme}
\end{figure}

\begin{definition}
    \subsection{Estremo superiore e Estremo inferiore}
    Sia \(A  \subset R \) un intervallo, chiuso o aperto .
    \begin{itemize}
        \item \(\exists\ supA \in A\) \textbf{estremo superiore} di A se,
              \[\forall x \in A, \, x \leq supA.\]
              \[\forall r \in \mathbb{R}, \, \exists x \in A | x > r.\]
              è dunque il pi\'u piccolo dei maggioranti
        \item \(\exists\ infA \in A\) \textbf{estremo inferiore} di A se,
              \[\forall x \in A, \, x \geq infA.\]
              \[\forall r \in \mathbb{R}, \, \exists x \in A | x < r.\]
              è dunque il pi\'u grande dei minoranti
    \end{itemize}
\end{definition}
\begin{figure}[H]
    \centering
    \begin{tikzpicture}
        \draw[->] (-1,0) -- (5,0) ; % Linea numerica
        \foreach \x in  {1,3} % Punti
        \draw[shift={(\x,0)},color=black] (0pt,2pt) -- (0pt,-2pt);
        \filldraw [myblue] (1,0) circle (2pt); % Cerchio pieno per a
        \draw [myblue] (3,0) circle (2pt); % Cerchio vuoto per b
        \filldraw [black] (1,0) circle (1pt); % Cerchio pieno per a
        \filldraw [black] (3,0) circle (1pt); % Cerchio vuoto per b
        \draw[thick, myblue](1,0) -- (3,0);
        % Aggiungi le parentesi e le etichette
        \node at (1,0.5) [above] {$infA$};
        \node at (3,0.5) [above] {$supA$};
        \node at (1,-0.5) [below] {$a$};
        \node at (3,-0.5) [below] {$b$};
        \node at (5,-0.5) [below] {\([a, b)\)};
    \end{tikzpicture}
    \caption{Estremo inferiore ovvero massimo dei minoranti, estremo superiore ovvero minimo dei maggioranti}
\end{figure}

\begin{definition}
    \subsection{Massimo e Minimo}
    Sia \(A \subset R, A = [a,b]\) un intervallo chiuso.
    \begin{itemize}
        \item \(x_M=\max{A} \in A\) è un \textbf{massimo} di A se,
              \[\forall x \in A, \, x \leq x_M.\]
        \item \(x_m=\min{A} \in A\) è un \textbf{minimo} di A se,
              \[\forall x \in A, \, x \geq x_m.\]
    \end{itemize}
\end{definition}

\begin{figure}[H]
    \centering
    \begin{tikzpicture}
        \draw[->] (-1,0) -- (5,0) ; % Linea numerica
        \filldraw [myblue] (1,0) circle (2pt); % Cerchio pieno per a
        \filldraw [myblue] (3,0) circle (2pt); % Cerchio vuoto per b
        \filldraw [black] (1,0) circle (1pt); % Cerchio pieno per a
        \filldraw [black] (3,0) circle (1pt); % Cerchio vuoto per b
        \draw[thick, myblue](1,0) -- (3,0);
        % Aggiungi le parentesi e le etichette
        \node at (1,0.5) [above] {$x_m=\min{A}$};
        \node at (3,0.5) [above] {$x_M=\max{A}$};
        \node at (1,-0.5) [below] {$a$};
        \node at (3,-0.5) [below] {$b$};
        \node at (5,-0.5) [below] {\([a, b]\)};
    \end{tikzpicture}
    \caption{Massimo e minimo}
\end{figure}

ma \textbf{non}
\begin{figure}[H]
    \centering
    \begin{tikzpicture}
        \draw[->] (-1,0) -- (5,0) ; % Linea numerica
        \draw [myblue] (1,0) circle (2pt); % Cerchio pieno per a
        \draw [myblue] (3,0) circle (2pt); % Cerchio vuoto per b
        % Croce al punto (1,0)
        \draw [black] (0.9,0.1) -- (1.1,-0.1);
        \draw [black] (0.9,-0.1) -- (1.1,0.1);
        % Croce al punto (3,0)
        \draw [black] (2.9,0.1) -- (3.1,-0.1);
        \draw [black] (2.9,-0.1) -- (3.1,0.1);
        \draw[thick, myblue](1,0) -- (3,0);
        % Aggiungi le parentesi e le etichette
        \node at (1,0.5) [above left] {$\nexists x_m=\min{A}$};
        \node at (3,0.5) [above right] {$\nexists x_M=\max{A}$};
        \node at (1,-0.5) [below] {$a$};
        \node at (3,-0.5) [below] {$b$};
        \node at (5,-0.5) [below] {\((a, b)\)};
    \end{tikzpicture}
    \caption{Massimo e minimo $\nexists$ in intervalli aperti}
\end{figure}

\section{Fattoriali e Coefficienti binomiali}
\begin{definition}
    \subsection{Fattoriale}
    La definizione del fattoriale di un numero naturale \( n \) è data dalla seguente formula:
    \[n! = \begin{cases}
            1               & \text{se } n = 0, \\
            n \times (n-1)! & \text{se } n > 0.
        \end{cases}\]
    Questo significa che il fattoriale di un numero \( n \) è il prodotto di tutti i numeri interi positivi fino a \( n \). Per esempio, \[ 4! = 4 \times 3 \times 2 \times 1 = 24 \]
    \(n!\) rappresenta il numero di possibili \textbf{disposizioni di n oggetti distinti in sequenza}, o - che è lo stesso - il numero di possibili \textbf{permutazioni} di \(n\) oggetti ordinati
\end{definition}

\begin{figure}
    \centering
    \begin{tikzpicture}
        \begin{axis}[
                xlabel={\( n \)},
                ylabel={\( n! \)},
                xmin=0, xmax=4,
                ymin=0, ymax=24,
                grid=both,
                axis lines=middle,
                axis line style=->,
                grid style={line width=.1pt, draw=gray!10},
                major grid style={line width=.2pt,draw=gray!50},
                minor tick num=5,
                enlargelimits={abs=0.5},
                ticklabel style={font=\small},
                xlabel style={at={(ticklabel* cs:1)}, anchor=north west},
                ylabel style={at={(ticklabel* cs:1)}, anchor=south west},
                legend pos=south east
            ]
            \addplot+[myblue, thick, smooth] plot coordinates {
                    (0,1)
                    (1,1)
                    (2,2)
                    (3,6)
                    (4,24)
                    (5,120)
                };
            \legend{\(n!\)}
        \end{axis}
    \end{tikzpicture}
    \caption{Grafico del Fattoriale}
\end{figure}
\noindent\begin{minipage}{\textwidth}
    \begin{curiosity}
        \subsubsection{Perché (0 ! = 1)}
        La definizione del fattoriale di un numero naturale \( n \), denotato come \( n! \), è il prodotto di tutti i numeri interi positivi fino a \( n \). Tuttavia, questa definizione solleva una domanda: qual è il valore di \( 0! \)?
        Per convenzione matematica, il valore di \( 0! \) è definito come 1. Questa scelta può sembrare arbitraria, ma ci sono diverse ragioni matematiche che giustificano questa definizione:
        \begin{enumerate}
            \item \textbf{Consistenza con la definizione di fattoriale}: Se consideriamo il fattoriale di un numero come il prodotto di tutti i numeri interi fino a quel numero, allora \( 0! \) sarebbe il prodotto di nessun numero. In matematica, il prodotto di nessun fattore è definito come l'elemento neutro del prodotto, che è 1.
            \item \textbf{Coerenza con la formula del coefficiente binomiale}: In combinatoria, il coefficiente binomiale \( \binom{n}{k} \) rappresenta il numero di modi in cui si possono scegliere \( k \) elementi da un insieme di \( n \) elementi. La formula per il coefficiente binomiale è \( \binom{n}{k} = \frac{n!}{k!(n-k)!} \). Se si sceglie \( k = 0 \), cioè nessun elemento da un insieme di \( n \) elementi, dovrebbe esserci esattamente 1 modo per farlo, il che richiede che \( 0! = 1 \).
            \item \textbf{Convenienza nelle serie matematiche}: In molte serie matematiche e in analisi, avere \( 0! = 1 \) semplifica notevolmente le formule e i calcoli, evitando la necessità di trattare separatamente il caso in cui \( n = 0 \).
        \end{enumerate}
        Inoltre, consideriamo un esempio operativo della definizione a ritroso dei valori del fattoriale:
        \[4! = 4 \times 3 \times 2 \times 1 = 24\]
        Da questo esempio, possiamo vedere come il fattoriale di un numero si costruisce moltiplicando tutti i numeri interi positivi fino a quel numero. Per definizione, possiamo poi vedere che:
        \[3! = \frac{4!}{4} = \frac{24}{4} = 6\]
        E proseguendo in questo modo:
        \[2! = \frac{3!}{3} = \frac{6}{3} = 2\]
        E infine:
        \[1! = \frac{2!}{2} = \frac{2}{2} = 1\]
        E da qui segue che:
        \[0! = \frac{1!}{1} = \frac{1}{1} = 1\]
        La definizione di \(0! = 1\) è quindi non solo una convenzione, ma anche una conseguenza logica della definizione generale di fattoriale.
    \end{curiosity}
\end{minipage}

\begin{curiosity}
    \subsubsection{Perché (a elevato alla 0 = 1)}
    Vista la curiosità precedente aggiungo anche questo esempio operativo per completezza, prendo un a = 4
    \[ \frac{4^a}{4^b} = 4^{a-b} \Longleftrightarrow 4^{a-b} \ = \frac{4^a}{4^b}  \]
    \[4^0 = 4^{1-1} = \frac{4^1}{4^1} = \frac{4}{4} \]
\end{curiosity}

\section{Scomposizione con Ruffini}
\subsection{How to Ruffini}
\begin{example}
    Vogliamo scomporre il polinomio
    \[2x^5 + x^3 + 5x - 8\]
    \begin{enumerate}
        \item Troviamo una radice \(r\) del polinomio, ovvero un valore di \(x\) per cui il polinomio si annulla.
        \item per farlo possiamo usare il teorema di Ruffini, che ci dice che se \(r\) è una radice del polinomio allora \(r\) è un divisore del termine noto, in questo caso \(r\) è un divisore di \(8\), quindi possibili valori di \(r\) sono \(\pm 1, \pm 2, \pm 4, \pm 8\).
        \item Proviamo con \(r = 1\), ottendendo
              \begin{align*}
                  2x^5 + x^3 + 5x - 8       & = 0 \\
                  2(1)^5 + (1)^3 + 5(1) - 8 & = 0 \\
                  2 + 1 + 5 - 8             & = 0 \\
                  0                         & = 0
              \end{align*}
        \item Quindi \(r = 1\) è una radice del polinomio.
        \item Ora dobbiamo scomporre il polinomio in fattori, per farlo usiamo le regola di Ruffini.
        \item Scriviamo i coefficienti del polinomio in una tabella, in questo caso abbiamo \(2, 0, 1, 5, -8\).
              \[
                  \ruffini{2,0,0,1,5,-8}{1}{ }{ }
              \]
              \[
                  \ruffini{2,0,0,1,5,-8}{1}{}{2}
              \]
              \[
                  \ruffini{2,0,0,1,5,-8}{1}{2,2}{2,2,2}
              \]
              \[
                  \ruffini{2,0,0,1,5,-8}{1}{2,2,2,3,8}{2,2,2,3,8,0}
              \]
        \item il polinomio ottenuto sarà
              \begin{align*}
                  2x^5 + x^3 + 5x - 8 = (x-1)(2x^4 + 2x^3 + 2x^2 + 3x + 8)
              \end{align*}
    \end{enumerate}
\end{example}
\end{document}
